% GENERATED FILE. Edit canonical lessons, then run npm run generate:books.
% canonical-source-hash: fnv1a64:b6159c622b08f794
% canonical-lessons: JA-C136-kazoku, JA-W136-zo, JA-C136-kazoeru, JA-C136-mizu, JA-W136-zu, JA-C136-suzushii, JA-C136-kaze, JA-W136-ze, JA-R136-family-and-water

\chapter{Family and Water: Zo, Zu and Ze}
\label{ch:ja-family-and-water-zo-zu-ze}
\hlchaptermodality{136}

\begin{chapteropening}
\textbf{By the end of this chapter:} \emph{I can write zo, zu and ze by adding the voicing mark to signs I already write, and read words that use them, such as kazoku, mizu and kaze.}

Read kazoku, kazoeru, mizu, suzushii and kaze, write the three new signs beside the signs they voice, write kazoku and kazoeru from memory and say which one is a verb, and write kaze.
\end{chapteropening}

\section[\texorpdfstring{\ja{かぞく} (kazoku)}{かぞく (kazoku)}]{\ja{かぞく} (kazoku) — a family; my family}

\label{lesson:JA-C136-kazoku}



\begin{quote}
Three recalls before the new word.

\begin{itemize}
\raggedright
  \item \emph{Write it:} \textbf{\ja{げ}} from memory — \textbf{R2}, five lessons back
  \item \emph{Write it:} \textbf{\ja{ば}} from memory — \textbf{R3}, twenty lessons back
  \item \emph{Recall:} say \emph{drowsy} — \textbf{R4}, eighty lessons back
\end{itemize}
\end{quote}

\subsection*{You'll want to know: \ja{かぞく}}
\textbf{\ja{かぞく}} — \emph{kazoku} — \textquotedblleft{}a family\textquotedblright{}, the people you live with or grew up with. Said of your own, it means \emph{my family}.

Three beats: \emph{ka–zo–ku}. \textbf{\ja{か}} and \textbf{\ja{く}} you already write. The middle sign is \textbf{\ja{そ}} with the two short strokes of the voicing mark at its upper right, so \emph{so} becomes \emph{zo}. The next lesson writes it.

\subsection*{Your turn}
\begin{itemize}
\raggedright
  \item \emph{Say it:} \emph{kazoku}
  \item \emph{Say it:} \emph{kazoku}, then count its beats aloud — three
  \item \emph{Recall:} say which sign in \textbf{\ja{かぞく}} carries the two-stroke mark, and name the sign under it
\end{itemize}

\begin{tcolorbox}[breakable,colback=teal!4,colframe=teal!35!black,title={Before you move on}]
What is \textquotedblleft{}a family\textquotedblright{}? (\textbf{\ja{かぞく}}, \emph{kazoku}.) How many beats? (\textbf{Three}.)
\end{tcolorbox}
\section[\texorpdfstring{\ja{ぞ} (zo)}{ぞ (zo)}]{\ja{ぞ} — \ja{そ}, voiced to zo}

\label{lesson:JA-W136-zo}



\begin{quote}
Four recalls, then the sign.

\begin{itemize}
\raggedright
  \item \emph{Recall:} say \emph{a family} — \textbf{R1}, one lesson back
  \item \emph{Recall:} say \emph{well}, as in healthy — \textbf{R2}, five lessons back
  \item \emph{Recall:} say \emph{a place} — \textbf{R3}, twenty lessons back
  \item \emph{Recall:} say \emph{relieved} — \textbf{R4}, eighty lessons back
\end{itemize}
\end{quote}

\subsection*{Script — the change}
\hlblockfigure{\detokenize{figures/JA-W136-zo-filmstrip.pdf}}{How \ja{ぞ} is written, stroke by stroke}

\begin{quote}
\textbf{\ja{そ}} \emph{so} + \textbf{\ja{゛}} $\to$ \textbf{\ja{ぞ}} \emph{zo}
\end{quote}

Write \textbf{\ja{そ}} in full: one stroke that zigzags twice before it drops into an open curve. Then add the two short strokes of the voicing mark at the upper right, the left one first. The base shape does not change.

The mark voices the opening consonant: \emph{s} becomes \emph{z}, and the vowel stays \emph{o}. It is the change you made on \textbf{\ja{さ}} for \textbf{\ja{ざ}}, in \textbf{\ja{ございます}}.

With \textbf{\ja{か}} in front and \textbf{\ja{く}} after it, you can now write the word from the last lesson: \textbf{\ja{かぞく}}, \emph{kazoku}, a family.

\subsection*{Your turn}
\begin{enumerate}
\raggedright
  \item Write \textbf{\ja{そ}}, then \textbf{\ja{ぞ}}, and say \emph{so}, \emph{zo}.
  \item Write \textbf{\ja{さ} \ja{ざ}} and \textbf{\ja{そ} \ja{ぞ}} in pairs, and say each one.
  \item Write \textbf{\ja{かぞく}} from memory and say \emph{kazoku}.
\end{enumerate}

\begin{tcolorbox}[breakable,colback=teal!4,colframe=teal!35!black,title={Before you move on}]
Stop after one accurate form.

Source: \href{https://www.unicode.org/charts/PDF/U3040.pdf}{Unicode Hiragana chart}.
\end{tcolorbox}
\section[\texorpdfstring{\ja{かぞえる} (kazoeru)}{かぞえる (kazoeru)}]{\ja{かぞえる} (kazoeru) — to count}

\label{lesson:JA-C136-kazoeru}



\begin{quote}
Four recalls before the new word.

\begin{itemize}
\raggedright
  \item \emph{Write it:} \textbf{\ja{ぞ}} from memory — \textbf{R1}, one lesson back
  \item \emph{Recall:} say \emph{a bank} — \textbf{R2}, five lessons back
  \item \emph{Recall:} say \emph{to eat} — \textbf{R3}, twenty lessons back
  \item \emph{Recall:} say \emph{ordinary} — \textbf{R4}, eighty lessons back
\end{itemize}
\end{quote}

\subsection*{You'll want to know: \ja{かぞえる}}
\textbf{\ja{かぞえる}} — \emph{kazoeru} — \textquotedblleft{}to count\textquotedblright{}: one, two, three, out loud or on your fingers. This is the dictionary form, the one a word list gives. Said politely, it is \emph{kazoemasu}.

The new sign is spent at once. \textbf{\ja{か}}, \textbf{\ja{ぞ}}, \textbf{\ja{え}} and \textbf{\ja{る}}: four beats, \emph{ka–zo–e–ru}. It ends in \textbf{\ja{る}}, like the verb \textbf{\ja{たべる}}. You write every sign in it now.

\subsection*{Your turn}
\begin{itemize}
\raggedright
  \item \emph{Say it:} \emph{kazoeru}
  \item \emph{Say it:} \emph{kazoku}, then \emph{kazoeru}, and listen for the \emph{kazo} in both
  \item \emph{Write it:} \textbf{\ja{かぞえる}} from memory, four signs for four beats
\end{itemize}

\begin{tcolorbox}[breakable,colback=teal!4,colframe=teal!35!black,title={Before you move on}]
What does \ja{かぞえる} mean? (\textquotedblleft{}To count\textquotedblright{}.) Say it once more.
\end{tcolorbox}
\section[\texorpdfstring{\ja{みず} (mizu)}{みず (mizu)}]{\ja{みず} (mizu) — water; cold water}

\label{lesson:JA-C136-mizu}



\begin{quote}
Four recalls before the new word.

\begin{itemize}
\raggedright
  \item \emph{Recall:} say \emph{to count} — \textbf{R1}, one lesson back
  \item \emph{Write it:} \textbf{\ja{ぎ}} from memory — \textbf{R2}, five lessons back
  \item \emph{Write it:} \textbf{\ja{べ}} from memory — \textbf{R3}, twenty lessons back
  \item \emph{Recall:} say \emph{tasty}, or \emph{skilful} — \textbf{R4}, eighty lessons back
\end{itemize}
\end{quote}

\subsection*{You'll want to know: \ja{みず}}
\textbf{\ja{みず}} — \emph{mizu} — \textquotedblleft{}water\textquotedblright{}. It means cold or cool water: the water you drink, or ask for in a restaurant. Hot water has a word of its own.

Two beats: \emph{mi–zu}. \textbf{\ja{み}} you already write. The second sign is \textbf{\ja{す}} with the voicing mark, so \emph{su} becomes \emph{zu}. The next lesson writes it.

\subsection*{Your turn}
\begin{itemize}
\raggedright
  \item \emph{Say it:} \emph{mizu}
  \item \emph{Say it:} \emph{mizu}, then count its beats aloud — two
  \item \emph{Recall:} say which sign in \textbf{\ja{みず}} carries the two-stroke mark, and name the sign under it
\end{itemize}

\begin{tcolorbox}[breakable,colback=teal!4,colframe=teal!35!black,title={Before you move on}]
What is \textquotedblleft{}water\textquotedblright{}? (\textbf{\ja{みず}}, \emph{mizu}.) How many beats? (\textbf{Two}.)
\end{tcolorbox}
\section[\texorpdfstring{\ja{ず} (zu)}{ず (zu)}]{\ja{ず} — \ja{す}, voiced to zu}

\label{lesson:JA-W136-zu}



\begin{quote}
Three recalls, then the sign.

\begin{itemize}
\raggedright
  \item \emph{Recall:} say \emph{water} — \textbf{R1}, one lesson back
  \item \emph{Recall:} say \emph{study} — \textbf{R3}, twenty lessons back
  \item \emph{Recall:} say \emph{many}, \emph{a lot} — \textbf{R4}, eighty lessons back
\end{itemize}
\end{quote}

\subsection*{Script — the change}
\hlblockfigure{\detokenize{figures/JA-W136-zu-filmstrip.pdf}}{How \ja{ず} is written, stroke by stroke}

\begin{quote}
\textbf{\ja{す}} \emph{su} + \textbf{\ja{゛}} $\to$ \textbf{\ja{ず}} \emph{zu}
\end{quote}

Write \textbf{\ja{す}} in full: the horizontal, then the vertical with its loop. Then add the two short strokes of the voicing mark at the upper right, the left one first, clear of the horizontal.

As with \textbf{\ja{さ}} and \textbf{\ja{ざ}}, and \textbf{\ja{そ}} and \textbf{\ja{ぞ}}, the \emph{s} sign voices to \emph{z}: \emph{su} becomes \emph{zu}.

With \textbf{\ja{み}} in front of it, you can now write the word from the last lesson: \textbf{\ja{みず}}, \emph{mizu}, water.

\subsection*{Your turn}
\begin{enumerate}
\raggedright
  \item Write \textbf{\ja{す}}, then \textbf{\ja{ず}}, and say \emph{su}, \emph{zu}.
  \item Write \textbf{\ja{さ} \ja{ざ}}, \textbf{\ja{し} \ja{じ}} and \textbf{\ja{す} \ja{ず}} in pairs, and say each one.
  \item Write \textbf{\ja{みず}} from memory and say \emph{mizu}.
\end{enumerate}

\begin{tcolorbox}[breakable,colback=teal!4,colframe=teal!35!black,title={Before you move on}]
Stop after one accurate form.

Source: \href{https://www.unicode.org/charts/PDF/U3040.pdf}{Unicode Hiragana chart}.
\end{tcolorbox}
\section[\texorpdfstring{\ja{すずしい} (suzushii)}{すずしい (suzushii)}]{\ja{すずしい} (suzushii) — cool, pleasantly cool}

\label{lesson:JA-C136-suzushii}



\begin{quote}
Four recalls before the new word.

\begin{itemize}
\raggedright
  \item \emph{Recall:} write \textbf{\ja{ず}} — \textbf{R1}, one lesson back
  \item \emph{Recall:} say \emph{a family} — \textbf{R2}, five lessons back
  \item \emph{Recall:} say \emph{a newspaper} — \textbf{R3}, twenty lessons back
  \item \emph{Recall:} say \emph{few}, \emph{little} — \textbf{R4}, eighty lessons back
\end{itemize}
\end{quote}

\subsection*{You'll want to know: \ja{すずしい}}
\textbf{\ja{すずしい}} — \emph{suzushii} — \textquotedblleft{}cool\textquotedblright{}: the air on a summer evening, or a shady room on a hot day. It is a pleasant word; it is not \emph{cold}.

The new sign is spent at once. \textbf{\ja{す}}, \textbf{\ja{ず}}, \textbf{\ja{し}} and \textbf{\ja{い}}: four beats, \emph{su–zu–shi–i}. The plain \textbf{\ja{す}} and the voiced \textbf{\ja{ず}} stand side by side, so you can see the mark do its work. You write every sign in it now.

\subsection*{Your turn}
\begin{itemize}
\raggedright
  \item \emph{Say it:} \emph{suzushii}
  \item \emph{Say it:} \emph{mizu}, then \emph{suzushii}, and listen for the \emph{zu} in both
  \item \emph{Recall:} write \textbf{\ja{すずしい}}, with \textbf{\ja{す}} and then \textbf{\ja{ず}}
\end{itemize}

\begin{tcolorbox}[breakable,colback=teal!4,colframe=teal!35!black,title={Before you move on}]
What does \ja{すずしい} mean? (\textquotedblleft{}Cool\textquotedblright{}, pleasantly so.) Say it once more.
\end{tcolorbox}
\section[\texorpdfstring{\ja{かぜ} (kaze)}{かぜ (kaze)}]{\ja{かぜ} (kaze) — a wind; a cold}

\label{lesson:JA-C136-kaze}



\begin{quote}
Four recalls before the new word.

\begin{itemize}
\raggedright
  \item \emph{Recall:} say \emph{cool}, of a summer evening — \textbf{R1}, one lesson back
  \item \emph{Write it:} \textbf{\ja{ぞ}} from memory — \textbf{R2}, five lessons back
  \item \emph{Write it:} \textbf{\ja{ぶ}} from memory — \textbf{R3}, twenty lessons back
  \item \emph{Recall:} say \emph{refreshing} — \textbf{R4}, eighty lessons back
\end{itemize}
\end{quote}

\subsection*{You'll want to know: \ja{かぜ}}
\textbf{\ja{かぜ}} — \emph{kaze} — \textquotedblleft{}a wind\textquotedblright{}, the air that moves. The same two signs are also the word for \textquotedblleft{}a cold\textquotedblright{}, the illness you catch in winter. Kanji tell the two apart; in kana they look the same.

Two beats: \emph{ka–ze}. \textbf{\ja{か}} you already write. The second sign is \textbf{\ja{せ}} with the voicing mark, so \emph{se} becomes \emph{ze}. The next lesson writes it.

\subsection*{Your turn}
\begin{itemize}
\raggedright
  \item \emph{Say it:} \emph{kaze}
  \item \emph{Say it:} \emph{kaze}, then count its beats aloud — two
  \item \emph{Recall:} say which sign in \textbf{\ja{かぜ}} carries the two-stroke mark, and name the sign under it
\end{itemize}

\begin{tcolorbox}[breakable,colback=teal!4,colframe=teal!35!black,title={Before you move on}]
What is \textquotedblleft{}a wind\textquotedblright{}? (\textbf{\ja{かぜ}}, \emph{kaze}.) And \textquotedblleft{}a cold\textquotedblright{}? (The same: \textbf{\ja{かぜ}}.)
\end{tcolorbox}
\section[\texorpdfstring{\ja{ぜ} (ze)}{ぜ (ze)}]{\ja{ぜ} — \ja{せ}, voiced to ze}

\label{lesson:JA-W136-ze}



\begin{quote}
Three recalls, then the sign.

\begin{itemize}
\raggedright
  \item \emph{Recall:} say \emph{a wind}, or \emph{a cold} — \textbf{R1}, one lesson back
  \item \emph{Recall:} say \emph{to count} — \textbf{R2}, five lessons back
  \item \emph{Recall:} say \emph{calm} — \textbf{R4}, eighty lessons back
\end{itemize}
\end{quote}

\subsection*{Script — the change}
\hlblockfigure{\detokenize{figures/JA-W136-ze-filmstrip.pdf}}{How \ja{ぜ} is written, stroke by stroke}

\begin{quote}
\textbf{\ja{せ}} \emph{se} + \textbf{\ja{゛}} $\to$ \textbf{\ja{ぜ}} \emph{ze}
\end{quote}

Write \textbf{\ja{せ}} in full, in the order you learned it. Then add the two short strokes of the voicing mark at the upper right, the left one first, beside the top of the right-hand vertical.

The \emph{s} sign voices to \emph{z}: \emph{se} becomes \emph{ze}. With it you write every sign of the \emph{z} row: \textbf{\ja{ざ} \ja{じ} \ja{ず} \ja{ぜ} \ja{ぞ}}.

With \textbf{\ja{か}} in front of it, you can now write the word from the last lesson: \textbf{\ja{かぜ}}, \emph{kaze}, a wind, or a cold.

\subsection*{Your turn}
\begin{enumerate}
\raggedright
  \item Write \textbf{\ja{せ}}, then \textbf{\ja{ぜ}}, and say \emph{se}, \emph{ze}.
  \item Write \textbf{\ja{ざ} \ja{じ} \ja{ず} \ja{ぜ} \ja{ぞ}} in order, and say each one.
  \item Write \textbf{\ja{かぜ}} from memory and say \emph{kaze}.
\end{enumerate}

\begin{tcolorbox}[breakable,colback=teal!4,colframe=teal!35!black,title={Before you move on}]
Stop after one accurate form.

Source: \href{https://www.unicode.org/charts/PDF/U3040.pdf}{Unicode Hiragana chart}.
\end{tcolorbox}
\section[(dialogue)]{Family and water — the z row, complete}

\label{lesson:JA-R136-family-and-water}



\begin{quote}
Four recalls, then the review.

\begin{itemize}
\raggedright
  \item \emph{Write it:} \textbf{\ja{ぜ}} from memory — \textbf{R1}, one lesson back
  \item \emph{Recall:} say \emph{water} — \textbf{R2}, five lessons back
  \item \emph{Recall:} say \emph{a hospital} — \textbf{R3}, twenty lessons back
  \item \emph{Recall:} say \emph{honest}, \emph{straightforward} — \textbf{R4}, eighty lessons back
\end{itemize}
\end{quote}

\subsection*{Your turn — read the five words}
For each one below, say what it means.

Read these aloud:

\begin{itemize}
\raggedright
  \item \textbf{\ja{かぞく}} — a family
  \item \textbf{\ja{かぞえる}} — to count
  \item \textbf{\ja{みず}} — water
  \item \textbf{\ja{すずしい}} — cool
  \item \textbf{\ja{かぜ}} — a wind; a cold
\end{itemize}

Each of these words holds one of the three signs this chapter wrote. With them, the book writes every sign of the \emph{z} row, \textbf{\ja{ざ} \ja{じ} \ja{ず} \ja{ぜ} \ja{ぞ}}.

\subsection*{Your turn — write}
\begin{enumerate}
\raggedright
  \item \emph{Write it:} \textbf{\ja{そ} \ja{ぞ}}, \textbf{\ja{す} \ja{ず}} and \textbf{\ja{せ} \ja{ぜ}}, in pairs
  \item \emph{Write it:} \textbf{\ja{かぞく}} and \textbf{\ja{かぞえる}} from memory, and say which one is a verb
  \item Say \emph{kaze}. \emph{Write it:} \textbf{\ja{かぜ}}
\end{enumerate}

\begin{tcolorbox}[breakable,colback=teal!4,colframe=teal!35!black,title={Before you move on}]
A family? (\textbf{\ja{かぞく}}, \emph{kazoku}.) To count? (\textbf{\ja{かぞえる}}, \emph{kazoeru}.) Water? (\textbf{\ja{みず}}, \emph{mizu}.) Cool? (\textbf{\ja{すずしい}}, \emph{suzushii}.) A wind, or a cold? (\textbf{\ja{かぜ}}, \emph{kaze}.)
\end{tcolorbox}
